% =====================================================================
\section{Apartamientos del modelo ideal}
% =====================================================================

\subsection{Modulación del ancho de base (efecto Early)}
\label{sec:early}

Al aumentar la tensión inversa $|V_{CB}|$, la ZCE de la juntura C--B se ensancha y una parte
de ese ensanchamiento invade la base: el ancho cuasineutro
\[
  W = W_B - x_{nEB}(V_{EB}) - x_{nCB}(V_{CB})
\]
\textbf{disminuye} \diapo{2.3}. Como en MAD
\[
  I_C \simeq qA\frac{D_B}{W}\,p_{B0}\,e^{V_{EB}/V_T}
  \qquad\left(\text{y también } \beta\simeq\frac{1}{\frac{D_E}{D_B}\frac{N_B}{N_E}\frac{W}{L_E}+\frac12\frac{W^2}{L_B^2}}\right),
\]
$I_C$ y $\beta$ \textbf{crecen} con $|V_{CB}|$ aunque $V_{EB}$ se mantenga fija. En las curvas de
salida eso se ve como una pendiente: las rectas de la región activa, prolongadas hacia la
izquierda, se cortan aproximadamente en un mismo punto del eje, $V_{EC}=-V_A$, donde $V_A$ es
la \textbf{tensión de Early} \diapo{2.4}.

\begin{figure}[H]
\centering
\begin{tikzpicture}
\begin{axis}[
  width=0.8\textwidth,height=6cm,
  xmin=-45,xmax=12,ymin=0,ymax=3.2,
  axis lines=middle,
  xlabel={$V_{EC}$},ylabel={$I_C$},
  xlabel style={at={(ticklabel* cs:1)},anchor=north west},
  ytick=\empty,xtick={-40},xticklabels={$-V_A$},
  every axis plot/.append style={line width=1.2pt},
]
\foreach \k in {0.6,1.2,1.8,2.4}{
  \edef\tmp{\noexpand\addplot[azul,domain=0.4:12] {\k*(1+x/40)};
            \noexpand\addplot[azul,dashed,thin,domain=-40:0.4] {\k*(1+x/40)};}
  \tmp
}
\fill[azul] (axis cs:-40,0) circle (2pt);
\node[font=\scriptsize,anchor=west,align=left] at (axis cs:-36,2.6) {pendiente $\dfrac{\partial I_C}{\partial V_{EC}}\simeq\dfrac{I_C}{V_A}=\dfrac{1}{r_o}$};
\end{axis}
\end{tikzpicture}
\caption{Efecto Early: las curvas de salida extrapoladas convergen en $-V_A$ (Sedra, fig.~6.18;
\diapo{2.4}). Valores típicos: $V_A=\SIrange{20}{200}{V}$.}
\label{fig:early}
\end{figure}

\begin{clave}[Efecto Early]
\[
  I_C \simeq I_S\,e^{V_{EB}/V_T}\left(1+\frac{V_{EC}}{V_A}\right),
  \qquad
  r_o = \left(\frac{\partial I_C}{\partial V_{EC}}\right)^{-1} = \frac{V_A}{I_C}
  \ \text{\enhojatexto},
  \quad
  V_A \approx \frac{W}{\left|\dd W/\dd V_{CB}\right|}.
\]
\end{clave}

La salida del transistor no es una fuente de corriente ideal: tiene una resistencia en paralelo
$r_o$, finita. Para tener $V_A$ grande conviene que la ZCE de C--B penetre poco en la base, es
decir \textbf{$N_C<N_B$} (la ZCE se extiende sobre todo hacia el colector) y una base no
demasiado angosta. Hay un compromiso con $\beta$: base angosta = $\beta$ alto pero $V_A$ bajo.

\begin{ejemplo}[tensión de Early]
Para el transistor del ejemplo ($N_B=\SI{e16}{}$, $N_C=\SI{e15}{cm^{-3}}$, $W=\SI{1}{\micro m}$ con
$V_{CB}=\SI{-5}{V}$). La parte de la ZCE de C--B que cae en la base es
\[
  x_{nCB} = \sqrt{\frac{2\epsilon_r\epsilon_0}{q}\,\frac{N_C}{N_B(N_B+N_C)}\,(V_{bi}-V_{CB})},
  \qquad V_{bi}=V_T\ln\frac{N_BN_C}{n_i^2}=\SI{0.656}{V}.
\]
Con $\epsilon_r=11.8$: $x_{nCB}(\SI{-5}{V})=\SI{0.259}{\micro m}$ y
$\dd x_{nCB}/\dd|V_{CB}|=x_{nCB}/[2(V_{bi}-V_{CB})]=\SI{0.0229}{\micro m/V}$, así que
\[
  V_A\approx\frac{\SI{1}{\micro m}}{\SI{0.0229}{\micro m/V}}\approx\SI{44}{V},
  \qquad
  r_o=\frac{V_A}{I_C}=\frac{\SI{44}{V}}{\SI{1.58}{mA}}\approx\SI{28}{k\ohm}.
\]
\end{ejemplo}

\begin{ojo}[\diapo{2.3}]
``Asumiendo $V_{CB}>5V_T$'': en un PNP en MAD la juntura C--B está en inversa, o sea
$V_{CB}<0$. La condición correcta es $V_{CB}<-5V_T$ (es decir, $|V_{CB}|\gtrsim\SI{0.13}{V}$),
que permite despreciar $e^{V_{CB}/V_T}$ frente a 1.
\end{ojo}

\subsection{Perforación (\emph{punch-through})}

Si $|V_{CB}|$ sigue aumentando, la ZCE de C--B puede llegar a tocar la de E--B: $W\to0$
\diapo{2.5--2.6}. A partir de ahí la tensión de colector actúa directamente sobre la barrera
de la juntura E--B, la baja, y el emisor inyecta una corriente enorme que ya no controla la
base. La corriente queda limitada sólo por el circuito externo; si la potencia excede el
área de operación segura (SOA) el daño es irreversible.

La tensión de perforación $V_{PT}$ se obtiene de $W_B - x_{nEB} - x_{nCB}(V_{PT}) = 0$.

\begin{ejemplo}[¿perforación o avalancha?]
En el transistor del ejemplo, $W_B = W + x_{nEB}+x_{nCB} = 1 + 0.155 + 0.259 = \SI{1.41}{\micro m}$
($x_{nEB}=\SI{0.155}{\micro m}$ con $V_{EB}=\SI{0.65}{V}$, $V_{bi,EB}=\SI{0.835}{V}$). La perforación
ocurre cuando $x_{nCB}=W_B-x_{nEB}=\SI{1.26}{\micro m}$:
\[
  V_{bi}-V_{PT} = \frac{(\SI{1.26e-4}{cm})^2}{\frac{2\epsilon_r\epsilon_0}{q}\frac{N_C}{N_B(N_B+N_C)}}
  \approx\SI{134}{V}
  \quad\Rightarrow\quad V_{CB,PT}\approx\SI{-133}{V}.
\]
Una juntura de Si con $N_C=\SI{e15}{cm^{-3}}$ tiene una ruptura por avalancha de unos
\SI{300}{V}, así que en este transistor la perforación ocurriría antes. Para evitarla se
dopa más la base o se la hace más ancha (a costa de $\beta$).
\end{ejemplo}

\subsection{Ruptura por avalancha}

La juntura C--B es la que soporta la tensión alta \diapo{2.7}. Si el campo es suficientemente
grande, los portadores que la cruzan generan pares por \textbf{ionización por impacto} y la
corriente se multiplica por
\[
  M = \frac{1}{1-\left(|V_{CB}|/BV_{CB0}\right)^n}\qquad (2<n<10),
\]
donde $BV_{CB0}$ es la tensión de ruptura de la juntura C--B con el emisor abierto.

\paragraph{Base común (emisor abierto o $I_E$ fijo).} $I_C=M(\adc I_E+I_{CB0})$: la corriente
sólo diverge cuando $M\to\infty$, es decir en $|V_{CB}|=BV_{CB0}$.

\paragraph{Emisor común con base abierta ($I_B=0$).} Ahora $I_E=I_C$, y
\[
  I_C = M(\adc I_C+I_{CB0})\quad\Rightarrow\quad I_C=\frac{M\,I_{CB0}}{1-M\adc},
\]
que diverge cuando $M\adc=1$, mucho antes de que $M\to\infty$:
\begin{clave}[Tensión de ruptura en emisor común]
\[
  M\adc=1 \;\Rightarrow\; 1-\left(\frac{BV_{CE0}}{BV_{CB0}}\right)^n=\adc
  \;\Rightarrow\;
  BV_{CE0} = BV_{CB0}\,(1-\adc)^{1/n} \simeq \frac{BV_{CB0}}{\sqrt[n]{\beta}} .
\]
\end{clave}

\begin{intuicion}[por qué la ruptura en emisor común es más baja]
Es el mismo mecanismo que hace $I_{CE0}=(\beta+1)I_{CB0}$: en un PNP, los electrones generados
por avalancha en la ZCE de C--B son arrastrados a la base; con la base abierta no tienen salida,
polarizan más a la juntura E--B y el emisor inyecta más huecos, que generan más pares\ldots{}
Es una realimentación positiva con ganancia $\beta$, que dispara la corriente con una
multiplicación pequeña ($M=1/\adc\approx1.01$).
\end{intuicion}

\begin{figure}[H]
\centering
\begin{tikzpicture}
\begin{axis}[
  width=0.75\textwidth,height=5.8cm,
  xmin=0,xmax=105,ymin=0,ymax=5,restrict y to domain=0:6,
  xlabel={$V_{EC}$ (V)},ylabel={$I_C$},
  axis lines=left,ytick=\empty,
  xtick={52,100},xticklabels={$BV_{CE0}$,$BV_{CB0}$},
  every axis plot/.append style={line width=1.1pt,azul,samples=120},
]
\foreach \k in {0.0,0.8,1.6,2.4}{
  \edef\tmp{\noexpand\addplot[domain=0.5:51.6] {(\k+0.08)*(1+x/150)/(1-(x/52)^7)};}
  \tmp
}
\draw[gris,dashed] (axis cs:52,0) -- (axis cs:52,5);
\draw[gris,dashed] (axis cs:100,0) -- (axis cs:100,5);
\node[font=\scriptsize,anchor=west] at (axis cs:2,0.35) {$I_B=0$};
\end{axis}
\end{tikzpicture}
\caption{Ruptura en emisor común (\diapo{2.8}, Pierret fig.~11.10). Todas las curvas se
disparan cerca de $BV_{CE0}$, bastante antes que la ruptura de la juntura aislada $BV_{CB0}$.
Con $\beta=100$ y $n\approx7$ resulta $BV_{CE0}\approx0.52\,BV_{CB0}$.}
\label{fig:avalancha}
\end{figure}

\begin{ojo}[``$BV_{CE0}=0.52\,BV_{CB0}$'' \diapo{2.8}]
No es una regla general: corresponde a un caso particular ($\beta\approx100$ y $n\approx7$,
porque $100^{-1/7}=0.52$). Con $n=4$ y el mismo $\beta$ daría $0.32$. La relación general es
$BV_{CE0}\simeq BV_{CB0}\,\beta^{-1/n}$.
\end{ojo}

\subsection{Fototransistor}

Si la luz incide sobre la ZCE de C--B, los pares fotogenerados son separados por el campo: los
huecos van al colector y los electrones a la base \diapo{2.7}. Esa fotocorriente $I_L$ cumple
el papel de una corriente de base (aunque el terminal B esté abierto) y el transistor la
amplifica:
\[
  I_C = (\beta+1)\,I_L .
\]
El fototransistor es mucho más sensible que un fotodiodo, pero también mucho más lento (la
ganancia se paga con ancho de banda, como en el fotoconductor de la unidad 6).

\subsection{Otras desviaciones}

\begin{itemize}
  \item \textbf{Recombinación en la ZCE de E--B}: degrada $\beta$ a corrientes bajas (Gummel).
  \item \textbf{Alta inyección y resistencia de base}: degradan $\beta$ a corrientes altas.
  \item \textbf{Base gradual}: en los transistores reales la base se dopa por difusión o
  implantación y $N_B$ no es uniforme: aparece un campo interno que \emph{acelera} a los
  minoritarios, mejora $\aT$ y la velocidad.
  \item \textbf{Geometría real}: el ``efecto transistor'' ocurre sólo en el pequeño volumen
  vertical debajo del emisor \diapo{2.23}; el resto agrega resistencias ($r_b$, $r_c$) y
  capacidades parásitas.
\end{itemize}
